\begingroup
\shorthandoff{"<>} % Previene errores con babel y caracteres especiales en \texttt
\setlength{\arrayrulewidth}{1pt}
\arrayrulecolor{vlepsbased}
\begin{longtable}{|p{15cm}|}
  
  % --- CABECERA INICIAL ---
  \hline
  \rowcolor{maincolor} 
  \multicolumn{1}{|l|}{\textcolor{white}{\textbf{Caso de Uso}}} \\
  \hline
  \endfirsthead
  
  % --- CABECERA EN PÁGINAS SIGUIENTES ---
  \multicolumn{1}{c}{{\tablename\ \thetable{} -- Continuación}} \\
  \hline
  \rowcolor{maincolor} 
  \multicolumn{1}{|l|}{\textcolor{white}{\textbf{Caso de Uso (Continuación)}}} \\
  \hline
  \endhead
  
  % --- PIE DE PÁGINA INTERMEDIO ---
  \hline \multicolumn{1}{|r|}{\textit{Continúa en la siguiente página...}} \\ \hline
  \endfoot
  
  % --- PIE FINAL (Caption abajo) ---
  \hline
  \caption{Caso de Uso CU3 — Configuración Global (Singleton) por el Administrador.}
  \label{TB:CU3_SINGLETON}
  \endlastfoot

  % Fila 1: ID y Actor
  \rowcolor{ulepsbased}
  \begin{tabularx}{\linewidth}{@{}p{7.5cm}p{7.5cm}@{}}
    \textbf{Identificador:} CU3 & \textbf{Actor principal:} Administrador
  \end{tabularx} \\ \hline

  % Fila 2: Nombre
  \textbf{Nombre:} Configuración global del sistema (patrón Singleton) \\ \hline

  % Fila 3: Descripción
  \rowcolor{ulepsbased}
  \textbf{Descripción:} \newline 
  El administrador accede al panel de Django Admin para configurar los parámetros globales del sistema: modelo LLM del tutor, modelo LLM de moderación, modo de moderación y umbral de palabras para la moderación de salida. \\ \hline

  % Fila 4: Precondiciones
  \textbf{Precondiciones:} \newline
  1. El administrador tiene acceso al panel de Django Admin. \newline
  2. El usuario tiene permisos de staff. \\ \hline

  % Fila 5: Flujo principal
  \rowcolor{ulepsbased}
  \textbf{Flujo principal:} \newline
  1. El administrador inicia sesión en el panel de Django Admin con credenciales de staff.\newline
  2. El sistema muestra el panel de administración con la sección ``Configuración del sistema''.\newline
  3. El administrador hace clic en ``Configuración del sistema''. El sistema redirige automáticamente al formulario de edición de la única fila existente. \newline
  4. El sistema muestra el formulario organizado en dos fieldsets:\newline
  --- \textbf{Moderación de contenido}: campo \texttt{moderation\_mode} (desplegable: Input + Output, Solo Input, Solo Output, Desactivada) y el campo \texttt{mod\_word\_window} (entero positivo) que determina cada cuántas palabras se realiza el análisis asíncrono.\newline
  --- \textbf{Modelos LLM (Ollama)}: campos \texttt{llm\_model} y \texttt{moderation\_model}. El sistema consulta a Ollama en tiempo real para poblar estos campos como listas desplegables dinámicas con los modelos disponibles. \newline
  5. El administrador modifica los valores deseados y pulsa ``Guardar''.\newline
  6. El sistema sobreescribe el registro principal y actualiza los parámetros en la base de datos, garantizando que solo exista una fila (Singleton). \newline
  7. El sistema muestra confirmación de guardado exitoso. \\ \hline

  % Fila 6: Flujos alternativos
  \textbf{Flujos alternativos:} \newline
  \textbf{FA-01:} Servicio Ollama inaccesible al cargar: \newline
  1. El sistema consulta a Ollama la lista de modelos disponibles. \newline
  2. El servicio Ollama no responde o da error de conexión. \newline
  3. El sistema muestra los campos de selección como campos de texto libre para no bloquear al administrador. \newline
  4. El administrador modifica los datos si es necesario y el guardado se realiza sin errores. \newline
  \textbf{FA-02:} Modelo previamente guardado borrado: \newline
  1. El sistema consulta a Ollama la lista de modelos y obtiene los resultados correctamente. \newline
  2. El sistema detecta que el modelo guardado en la base de datos ya no se encuentra en la lista de instalados. \newline
  3. El sistema añade dicho modelo al desplegable, lo deja preseleccionado e incluye la advertencia «$\triangle$ no disponible». \\ \hline

  % Fila 7: Postcondiciones
  \rowcolor{ulepsbased}
  \textbf{Postcondiciones:} \newline
  1. Las peticiones de chat siguientes utilizarán los nuevos parámetros. \\ \hline

\end{longtable}
\endgroup